\documentclass[11pt]{article}
\usepackage[utf8]{inputenc}
\usepackage[T1]{fontenc}
\usepackage{amsmath,amssymb}
\usepackage[margin=1.1in]{geometry}
\usepackage{authblk}
\usepackage[colorlinks=true,linkcolor=blue,citecolor=blue,urlcolor=blue]{hyperref}

\title{A Theory That Does Not Face the Measurement Problem}

\author[1]{Claes Johnson}
\affil[1]{KTH Royal Institute of Technology, SE-100 44 Stockholm, Sweden \\
\small\texttt{claesjohnson@gmail.com}}

\date{\today}

\begin{document}
\maketitle

\begin{abstract}
The measurement problem is not a mystery about nature but a contradiction internal to a particular
formalism, and it can be stated as one: the wave function is complete, it always evolves linearly,
and measurements have determinate outcomes --- three propositions that cannot all hold. Every
interpretation of quantum mechanics is a choice of which to abandon. We observe that the
contradiction is incurred by three decisions taken in 1926 and not forced by any experiment: that
the state is a function on $3N$-dimensional configuration space, that its square is a probability,
and that its evolution is linear. A model of non-overlapping electron charge densities in
three-dimensional space takes none of the three. Its state is a definite configuration of charge at
every instant, so there is nothing to be made determinate; the square of the wave function is
charge, not probability, so the Born rule is not derived or denied but absent; and the evolution is
nonlinear, because each density moves in the potential the others generate, so superpositions of
macroscopically distinct states are not solutions and no apparatus can inherit one. There is no
collapse, no Heisenberg cut, and no cat. We then state the price, because a problem dissolved by a
theory that no longer accounts for the phenomena is not a problem solved. The same two features
that remove the contradiction --- three-dimensional charge and nonlinear evolution --- are what
place Bell correlations outside the model's reach: it has no entangled states, which is why Gisin's
superluminal-signalling objection to nonlinear quantum mechanics cannot gain a foothold, and
equally why the correlations measured at spacelike separation are not addressed. Spectral
discreteness and interference it does account for. Single-event statistics it does not derive. The
claim is therefore narrow and checkable: the measurement problem is not solved here, it is not
incurred, and the bill for not incurring it is itemised.
\end{abstract}

\section{The problem, stated so it can be checked}\label{sec:problem}

The measurement problem is often presented as a mystery, which makes it impossible to say whether
any proposal has addressed it. It is better than that. Maudlin showed that it can be put as a
plain inconsistency among three propositions~\cite{Maudlin1995}:

\begin{enumerate}\itemsep3pt
\item[(A)] The wave function of a system is \emph{complete}: it specifies all the physical
properties the system has.
\item[(B)] The wave function \emph{always} evolves linearly, by the Schr\"odinger equation.
\item[(C)] Measurements have \emph{determinate} outcomes: a pointer ends up in one place.
\end{enumerate}

The three cannot all be true. Prepare a system in a state that (B) carries into a sum of two
branches, couple it linearly to an apparatus, and linearity carries the sum to the apparatus too:
the pointer ends in a sum of positions. By (A) there is nothing further to appeal to, so by (C)
something must break.

Every interpretation is a decision about which proposition to surrender. Copenhagen suspends (B)
at a moment it declines to specify, which is the Heisenberg cut. Objective-collapse models modify
(B) explicitly and accept the experimental consequences. De Broglie--Bohm denies (A), adding
particle positions underneath. Everett keeps (A) and (B) and gives up (C) as normally understood,
retaining both branches. Each is coherent and each pays. That so many mutually incompatible options
remain live, none selected by the equations, is the signature of a formalism that does not say what
exists.

\section{The contradiction is incurred, not discovered}\label{sec:incurred}

What deserves emphasis is that the trio is not a report of experimental findings. Each of the three
propositions carries a modelling decision inside it, and all three decisions were taken in 1926.

\paragraph{The arena.} For one electron, Schr\"odinger's $\psi$ is a field on ordinary space and
admits a direct physical reading. For $N$ electrons the formalism was extended to a single function
$\Psi(x_1,\dots,x_N)$ on $3N$-dimensional configuration space. That extension is a postulate, not a
consequence: a one-electron equation fixes nothing about the dimension of the space in which the
many-electron problem is posed. And it is what makes a ``branch'' available in the first place, for
a function on configuration space can superpose alternatives that are mutually exclusive in
ordinary space.

\paragraph{The reading.} $|\Psi|^2$ was interpreted as a probability density. For one electron the
same quantity reads equally well as a charge density --- electrons per unit volume, a physical
field. The probabilistic reading is an interpretation laid on the mathematics, not extracted from
it, and (C) only becomes problematic once the state is a catalogue of possibilities rather than a
distribution of stuff.

\paragraph{The linearity.} Linearity of the evolution is what transmits a superposition from system
to apparatus. It is also a strong structural assumption, and the one that makes the problem
contagious: without it, a sum of two solutions is not a solution, and the pathology has no vehicle.

None of the three was compelled by measurement. They were chosen because together they reproduced
the spectra, and they have been kept because the combination works where it can be computed.

\section{A three-dimensional charge model}\label{sec:realqm}

The alternative examined here keeps Coulomb interaction and the kinetic form of the Schr\"odinger
functional, and extends them to many electrons without leaving three-dimensional
space~\cite{RealQM}. Instead of one antisymmetric function on $\mathbb{R}^{3N}$, there are $N$
real-valued functions $\psi_i$ on $\mathbb{R}^3$, each supported on its own domain $\Omega_i$, the
domains non-overlapping:
\begin{equation}
\Omega_i \cap \Omega_j = \emptyset \quad (i \neq j), \qquad
\rho(x) = \sum_{i=1}^{N} \psi_i^2(x), \qquad x \in \mathbb{R}^3 .
\label{eq:state}
\end{equation}
Each $\psi_i^2$ is the unit charge density of electron $i$ in its territory; $\rho$ is a single
real field on physical space. Electron $i$ moves in the potential generated by the \emph{other}
densities, the domain boundaries are free and located by the energy, and exclusion is geometric
rather than a symmetry of a high-dimensional function. There is no spin and $Z$ is the only input.

Nothing here is claimed to be exact. It is a postulate, like the one it replaces, and it is to be
judged by what it predicts and by what it costs. The relevant point for the present question is
structural: with the state given by \eqref{eq:state}, the three propositions of
Section~\ref{sec:problem} read quite differently.

\subsection*{(A) Completeness}

The state \eqref{eq:state} is complete and nothing is hidden beneath it, so this is not where the
model departs. But completeness is no longer a worry, because what the state specifies is already
determinate: a charge density has a value at every point of ordinary space. The model is not a
hidden-variable theory --- it adds nothing underneath an amplitude --- because there is no
amplitude for anything to be underneath.

\subsection*{(B) Linearity, and its absence}

This is the decisive difference, and it is technical rather than philosophical. The evolution is
\emph{nonlinear}. Each density moves in the potential the others generate, and that potential is a
functional of the densities:
\begin{equation}
-\tfrac12\nabla^2\psi_i \;-\; K\psi_i \;+\; \Big(\textstyle\sum_{j \neq i} P[\psi_j^2]\Big)\psi_i
\;=\; \varepsilon_i \psi_i ,
\qquad -\nabla^2 P[\psi_j^2] = 4\pi \psi_j^2 ,
\label{eq:nonlinear}
\end{equation}
with the supports $\Omega_i$ themselves determined by the solution. The coupling through $P$ is
quadratic in the unknowns and the free-boundary condition is nonlinear in them.

The consequence is immediate: \emph{a sum of two solutions of \eqref{eq:nonlinear} is not a
solution}. Superposition is not a property of this model's state space. One may of course add two
charge densities, but the sum is simply another charge density --- a different distribution of
stuff, not a disjunction of alternatives awaiting resolution. So there is no mechanism by which a
system can place an apparatus into a sum of pointer positions, because the required intermediate
object is not a state of the theory at all.

This removes the engine of the contradiction rather than negotiating with its conclusion. The step
in Section~\ref{sec:problem} that carries the branch structure from system to apparatus is
precisely the step that requires (B), and (B) fails here.

\subsection*{(C) Determinate outcomes}

Outcomes are determinate because states were determinate before, during, and after any
interaction. Apparatus and system are described by the same equations and the same kind of object,
so a ``measurement'' is an interaction like any other, with no special status and no distinguished
moment. The theory therefore needs no rule about when measurement occurs, which is the rule
standard quantum mechanics cannot supply from within its own dynamics.

\section{What is no longer needed}\label{sec:gone}

Three familiar items disappear, and it is worth being explicit that they disappear rather than
being explained.

\paragraph{The Born rule is absent, not derived.} $\psi_i^2$ is charge. There is no probability
anywhere in \eqref{eq:state} or \eqref{eq:nonlinear} to be interpreted, so the question of why
probabilities go as the square does not arise. This is a stronger and also a narrower claim than
deriving the rule: the model dispenses with the rule because it dispenses with the ensemble reading,
and it must then account for observed statistics by other means or admit that it has not
(Section~\ref{sec:price}).

\paragraph{Collapse is absent.} There is one evolution law and it applies always. Nothing needs to
be converted from possibility into actuality, because the state was never a catalogue of
possibilities.

\paragraph{The cut is absent.} Since there is no second law to switch to, there is no boundary to
place between quantum system and classical observer, and no need to decide which side an apparatus,
a detector, a retina or an experimenter falls on. Schr\"odinger's cat does not arise: the cat is
charge in three dimensions, in one configuration, at all times. The thought experiment was a
reductio of treating an amplitude on configuration space as the reality, and there is no such
amplitude here to be reduced.

\section{What a dissolution owes}\label{sec:owes}

A theory can always avoid a problem by not addressing the phenomena that generate it, and that is
no achievement. So the honest test of the above is whether the model still accounts for what made
the measurement problem pressing. Two things it does.

\paragraph{Discreteness.} Sharp spectral lines are what originally forced quantisation into the
theory, and they come out here without a quantisation postulate: the stationary configurations of
\eqref{eq:nonlinear} form a discrete set because the boundary-value problem on a bounded domain
does, and the observed frequencies are differences of their energies. The relation $E = h\nu$ then
functions as a conversion between an atomic energy scale and a radiation frequency scale rather
than as a statement that energy comes in lumps~\cite{RealQM}. Discreteness is a property of the
solution set, which is where a wave theory should put it.

\paragraph{Interference.} The equation is a wave equation and interference is what wave equations
do. A charge density passing an obstacle with two openings redistributes itself into fringes by
the ordinary mechanism, in three-dimensional space, with no appeal to a particle choosing a slit or
to amplitudes for histories. What the model does \emph{not} supply here is the single-detection
event --- see below.

\section{The price, itemised}\label{sec:price}

Three debts, stated plainly, because the value of the argument in Section~\ref{sec:realqm} depends
entirely on these being declared rather than discovered by a reader.

\paragraph{Bell correlations are outside the model's reach, and for the same reason that the
measurement problem is.} This is the central honest point of the paper, and it is not a coincidence
of two unrelated limitations. Correlated outcomes at spacelike separation, measured since Aspect
and now beyond reasonable dispute, exclude local hidden-variable accounts. The model has no
entangled states: there is no configuration space in which a two-particle state could fail to
factorise, and \eqref{eq:state} describes bound charge sharing one spatial domain, not pairs
separated by laboratory distances with correlated results. So the correlations are not predicted,
not explained, and not reproduced --- they are simply not in the theory's range, which is the
position taken throughout this programme for spin and magnetism as well.

There is a sharp way to see that this is one fact rather than two. Gisin showed that a deterministic
nonlinear modification of quantum mechanics permits superluminal signalling, because nonlinear
evolution destroys the equivalence of different ensembles representing the same mixed state, and
the construction runs on \emph{entangled} pairs~\cite{Gisin1990}. The model here is nonlinear and
so would be exposed to exactly that argument --- except that the argument needs entangled states,
and there are none. The nonlinearity which dissolves the measurement problem is therefore shielded
from Gisin's objection by the very same absence that leaves Bell experiments unaddressed. One
cannot keep one side of this trade and refuse the other, and we do not claim to.

\paragraph{Single-event statistics are not derived.} A detector registers one click in one place,
and repeated preparations give frequencies that standard quantum mechanics predicts via the Born
rule. Having dispensed with that rule, this model owes an account of those frequencies and does not
currently have one. A deterministic field theory must locate the apparent randomness somewhere ---
in the preparation, in the apparatus, in sensitivity to initial conditions --- and saying which,
with a calculation, is unfinished work. This is the most serious of the three debts, because it is
the place where the measurement problem's original motivation partly survives translation.

\paragraph{The model's own accuracy is a few per cent.} It is not exact and does not claim to be;
its energies carry errors of a few per cent from one electron to ten thousand~\cite{RealQM}. The
present argument is about structure, not accuracy, and nothing above depends on the per-cent
figures --- but a reader should not take a dissolution of the measurement problem as evidence that
the model is quantitatively superior where standard methods can be computed. It is not.

\section{Conclusion}\label{sec:conclusion}

The measurement problem is a contradiction among three propositions, each of which encodes a choice
made when the one-electron wave equation was extended to many electrons. A model of non-overlapping
charge densities in three-dimensional space makes different choices and the contradiction does not
form: the state is determinate, so nothing must be made determinate; the square of the wave
function is charge, so no probability requires interpreting; and the evolution is nonlinear, so no
superposition propagates into an apparatus. Collapse, the cut and the cat are not resolved but
unneeded.

This is a modest claim and should be read as one. It is not that a hundred years of foundational
difficulty was unnecessary, nor that this model is the right physics --- its reach stops well short
of the experiments where the measurement problem is now most sharply posed. It is that the
difficulty is a feature of a formalism rather than of nature, which is demonstrated by exhibiting a
formalism for the same matter that does not have it, together with an explicit statement of what
that formalism cannot yet do. The interesting question is then not how to interpret quantum
mechanics but which of 1926's three decisions were forced --- and on the evidence of a computable
three-dimensional alternative, fewer than is usually supposed.

\begin{thebibliography}{9}

\bibitem{Maudlin1995} T.~Maudlin, \emph{Three measurement problems}, Topoi \textbf{14} (1995)
7--15.

\bibitem{Gisin1990} N.~Gisin, \emph{Weinberg's non-linear quantum mechanics and supraluminal
communications}, Phys.\ Lett.\ A \textbf{143} (1990) 1--2.

\bibitem{RealQM} C.~Johnson, \emph{Real Quantum Mechanics}, and the computational Gallery,
\url{https://claes542.github.io/RealMolecule/gallery.html}.

\end{thebibliography}

\end{document}
